%% Generated by tables/make_wood_tables.py -- do not edit.
\begin{deluxetable*}{lccccccccccc}
\tablecaption{Lexington \HII{} region benchmark HII40 (40~kK blackbody):
the published multi-code results of \citet[][their Table 1]{wood2004}, with a
\textsc{Cloudy} c25.00 calculation and the \MoCHII{} result appended.  Line
intensities are relative to H$\beta$.\label{tab:wood40}}
\tablewidth{0pt}
\tabletypesize{\footnotesize}
\tablehead{\colhead{Quantity} & \colhead{GF} & \colhead{C25} & \colhead{HN} & \colhead{DP} & \colhead{TK} & \colhead{PH} & \colhead{RS} & \colhead{RR} & \colhead{BE} & \colhead{WME} & \colhead{\MoCHII{}}}
\startdata
$L(\mathrm{H}\beta)/10^{37}\,\mathrm{erg\,s^{-1}}$ & 2.06 & 2.06 & 2.02 & 2.02 & 2.10 & 2.05 & 2.07 & 2.05 & 2.02 & 2.01 & 1.94 \\
\HeI{} 5876 & 0.119 & 0.100 & 0.112 & 0.113 & 0.116 & 0.118 & 0.116 & \nodata & 0.114 & 0.114 & 0.106 \\
C\,\textsc{ii}] 2325+ & 0.157 & 0.189 & 0.141 & 0.139 & 0.110 & 0.166 & 0.096 & 0.178 & 0.148 & 0.172 & 0.174 \\
C\,\textsc{iii}] 1907+1909 & 0.071 & 0.062 & 0.076 & 0.069 & 0.091 & 0.060 & 0.066 & 0.074 & 0.041 & 0.078 & 0.057 \\
$[$N\,\textsc{ii}$]$ 122\,$\mu$m & 0.027 & 0.030 & 0.037 & 0.034 & \nodata & 0.032 & 0.035 & 0.030 & 0.036 & 0.031 & 0.032 \\
$[$N\,\textsc{ii}$]$ 6584+6548 & 0.669 & 0.669 & 0.817 & 0.725 & 0.69 & 0.736 & 0.723 & 0.807 & 0.852 & 0.742 & 0.702 \\
$[$N\,\textsc{ii}$]$ 5755 & 0.0050 & 0.0060 & 0.0054 & 0.0050 & \nodata & 0.0064 & 0.0050 & 0.0068 & 0.0061 & 0.0057 & 0.0060 \\
$[$N\,\textsc{iii}$]$ 57.3\,$\mu$m & 0.306 & 0.286 & 0.261 & 0.311 & \nodata & 0.292 & 0.273 & 0.301 & 0.223 & 0.308 & 0.288 \\
$[$O\,\textsc{i}$]$ 6300+6363 & 0.0094 & 0.0112 & 0.0086 & 0.0088 & 0.012 & 0.0059 & 0.0070 & \nodata & 0.0065 & 0.011 & 0.0170 \\
$[$O\,\textsc{ii}$]$ 7320+7330 & 0.029 & 0.033 & 0.030 & 0.031 & 0.023 & 0.032 & 0.024 & 0.036 & 0.025 & 0.033 & 0.035 \\
$[$O\,\textsc{ii}$]$ 3726+3729 & 1.94 & 2.32 & 2.17 & 2.12 & 1.6 & 2.19 & 1.88 & 2.26 & 1.92 & 2.23 & 2.32 \\
$[$O\,\textsc{iii}$]$ 52+88\,$\mu$m & 2.35 & 2.40 & 2.10 & 2.26 & 2.17 & 2.34 & 2.29 & 2.34 & 2.28 & 2.42 & 2.43 \\
$[$O\,\textsc{iii}$]$ 5007+4959 & 2.21 & 2.46 & 2.38 & 2.20 & 3.27 & 1.93 & 2.17 & 2.08 & 1.64 & 2.34 & 2.37 \\
$[$O\,\textsc{iii}$]$ 4363 & 0.00235 & 0.0046 & 0.0046 & 0.0041 & 0.0070 & 0.0032 & 0.0040 & 0.0035 & 0.0022 & 0.0044 & 0.0044 \\
$[$Ne\,\textsc{ii}$]$ 12.8\,$\mu$m & 0.177 & 0.200 & 0.195 & 0.192 & \nodata & 0.181 & 0.217 & 0.196 & 0.212 & 0.192 & 0.206 \\
$[$Ne\,\textsc{iii}$]$ 15.5\,$\mu$m & 0.294 & 0.216 & 0.264 & 0.270 & 0.35 & 0.429 & 0.350 & 0.417 & 0.267 & 0.263 & 0.284 \\
$[$Ne\,\textsc{iii}$]$ 3869+3968 & 0.084 & 0.058 & 0.087 & 0.071 & 0.092 & 0.087 & 0.083 & 0.086 & 0.053 & 0.063 & 0.072 \\
$[$S\,\textsc{ii}$]$ 6716+6731 & 0.137 & 0.304 & 0.166 & 0.153 & 0.315 & 0.155 & 0.133 & 0.130 & 0.141 & 0.18 & 0.323 \\
$[$S\,\textsc{ii}$]$ 4068+4076 & 0.0093 & 0.0142 & 0.0090 & 0.0100 & 0.026 & 0.0070 & 0.005 & 0.0060 & 0.0060 & 0.0094 & 0.0140 \\
$[$S\,\textsc{iii}$]$ 18.7\,$\mu$m & 0.627 & 0.504 & 0.750 & 0.726 & 0.535 & 0.556 & 0.567 & 0.580 & 0.574 & 0.623 & 0.509 \\
$[$S\,\textsc{iii}$]$ 9532+9069 & 1.13 & 1.00 & 1.19 & 1.16 & 1.25 & 1.23 & 1.25 & 1.28 & 1.21 & 1.44 & 1.01 \\
$10^{3}\times\Delta(\mathrm{BC}\,3645)$/\AA & 4.88 & \nodata & \nodata & 4.95 & \nodata & 5.00 & 5.70 & \nodata & 5.47 & 4.96 & \nodata \\
$T_{\mathrm{inner}}$ / K & 7719 & 7812 & 7668 & 7663 & 8318 & 7440 & 7644 & 7399 & 7370 & 7743 & 7540 \\
$\langle T[N_{\mathrm{p}}N_{\mathrm{e}}]\rangle$ / K & 7940 & 8210 & 7936 & 8082 & 8199 & 8030 & 8022 & 8060 & 7720 & 8183 & 8202 \\
$R_{\mathrm{out}}/10^{19}$\,cm & 1.46 & 1.47 & 1.46 & 1.46 & 1.45 & 1.46 & 1.47 & 1.46 & 1.46 & 1.45 & 1.45 \\
$\langle\mathrm{He^{+}}\rangle/\langle\mathrm{H^{+}}\rangle$ & 0.787 & 0.741 & 0.727 & 0.754 & 0.77 & 0.764 & 0.804 & 0.829 & 0.715 & 0.771 & 0.704 \\
\enddata
\tablecomments{Codes, in the notation of \citet{wood2004}: GF =
Ferland's \textsc{Cloudy}; HN = Netzer's \textsc{Ion}; DP = P\'equignot's
\textsc{Nebu}; TK = Kallman's \textsc{XStar}; PH = Harrington's code;
RS = Sutherland's \textsc{Mappings}; RR = Rubin's \textsc{Nebula};
BE = Ercolano's \textsc{Mocassin}; WME = the code of \citet{wood2004}
itself.  Columns GF--BE agree entry by entry with Tables 4 and 5 of
\citet{ercolano2003}, the compilation of the published columns.
C25 is \textsc{Cloudy} c25.00 \citep{cloudy2025} run for this paper and
provides a uniform one-dimensional reference alongside the published
multi-code spread.
Its physics input is the unmodified c25.00 test-suite file
\code{tsuite/auto/hii\_paris.in} (HII40) or \code{hii\_coolstar.in} (HII20)
--- \textsc{Cloudy}'s own maintained version of these benchmarks --- with
only output commands appended; line luminosities come from a \code{save line
list absolute} file and are divided by the \textsc{Cloudy} H\,\textsc{i}
4861.32\,\AA{} luminosity, and the \textsc{Cloudy} blends summed for each row
are listed in \code{tables/make\_wood\_tables.py}.  Unlike the published
columns, C25 runs to the stopping criterion of its own deck rather than to a
prescribed outer radius --- an electron fraction of $10^{-2}$ at HII40, the
1000~K temperature floor at HII20 --- which mainly affects the
low-ionization rows fed by the partly neutral edge.
The \MoCHII{} column uses continuous stellar photon energies sampled from
the Planck spectrum with a Sobol sequence (seed 20260728), and evaluates
photoionization cross sections at each sampled energy; the 32-bin ionizing
array is a diagnostic tally only.  It uses the explicit Monte Carlo diffuse
field including the \HeI{} channels (case~A, matching the transfer treatment
of the published codes), and Badnell 2023 recombination with the Mao \&
Kaastra 2016 ground-level split, at $2^{25}$ stellar packets with the
volume-integrated convergence criterion; the ionizing luminosity is set from
the benchmark $L({\rm BB})$ so that $Q({\rm H})$ matches the specification;
blends are the sum of the
\MoCHII{} components listed in \code{tables/make\_wood\_tables.py}, and
\MoCHII{} wavelengths are vacuum values.  The \MoCHII{} structural entries
are computed from the same run: $T_{\mathrm{inner}}$ is the mean $\Te$ of the
cells at the illuminated inner edge, $R_{\mathrm{out}}$ the radius at which
the shell-averaged $x_{\mathrm{HII}}$ falls through 0.5, and the helium ratio
follows Eq.~(17) of \citet{ercolano2003},
$\langle\mathrm{He^{+}}\rangle/\langle\mathrm{H^{+}}\rangle =
\int \nelec n(\mathrm{He^{+}})\,{\rm d}V / \int \nelec n_{\mathrm{p}}\,{\rm d}V
\times n_{\mathrm{H}}/n_{\mathrm{He}}$.  Where the ionized gas still fills the
outermost shell of the prescribed sphere, $R_{\mathrm{out}}$ would be set by the
grid rather than predicted, and the entry is left empty.  The Balmer jump is
not computed.  The H$\beta$ row is headed $10^{36}$
erg\,s$^{-1}$ in the published table; the values listed there are the
$10^{37}$~erg\,s$^{-1}$ values of \citet{ercolano2003}, and the heading is
corrected accordingly here.}
\end{deluxetable*}
